\documentclass[11pt]{article}
\usepackage[margin=0.9in]{geometry}
\usepackage{graphicx,amsmath,amssymb,booktabs,hyperref,xcolor}
\graphicspath{{/tmp/cheby_h0/figs/three_way/}}
\hypersetup{colorlinks=true,linkcolor=blue,urlcolor=blue}
\setlength{\parskip}{6pt}
\title{Three tabulated-$\mu$ operators side-by-side:\\
       Chebyshev strict-$h{=}0$ vs Linear-CDF kernel vs Linear-CDF strict-$h{=}0$}
\author{Fixed-Point Factory --- $\gamma$-sweep, same architecture three ways}
\date{June 2026}

\begin{document}
\maketitle

\section*{The three operators}

All three share the tabulated-$\mu$ architecture (build $\mu(p, u_k)$
table, look up per cube cell, run CRRA clearing). They differ in:

\begin{center}
\begin{tabular}{lccc}
\toprule
& \textbf{Cheb strict $h{=}0$} & \textbf{Lin-CDF kernel} & \textbf{Lin-CDF strict $h{=}0$} \\
\midrule
$u$-grid                & Lobatto-atanh & Gaussian quantiles & Gaussian quantiles \\
$P$ representation      & tensor Cheb poly & piecewise bilinear & piecewise bilinear \\
contour integral        & strict $h{=}0$ line int. & Gaussian band $h{=}0.5$ & strict $h{=}0$ line int. \\
POU partition-of-unity  & yes & --- (kernel) & yes \\
ms / $\Phi$ (G=7)        & 540 & 6 & 7 \\
\bottomrule
\end{tabular}
\end{center}

The third operator (NEW): strict $h{=}0$ line integral, no kernel, just
finds piecewise-linear root crossings on the CDF grid, weighted by
partition-of-unity over both axes.

\section{Convergence floor per $\gamma$}

\begin{center}
\includegraphics[width=\linewidth]{02_floor_three.png}
\end{center}

\begin{center}
\begin{tabular}{lccc}
\toprule
$\gamma$ & Cheb strict $\|F\|_\infty$ & Lin-CDF kernel $\|F\|_\infty$ & Lin-CDF strict $\|F\|_\infty$ \\
\midrule
0.1 & $4.4\!\times\!10^{-16}\,\star$ & $2.2\!\times\!10^{-16}\,\star$ & $6.0\!\times\!10^{-2}$ \\
0.3 & $6.4\!\times\!10^{-4}$         & $2.2\!\times\!10^{-16}\,\star$ & $2.8\!\times\!10^{-2}$ \\
0.5 & $1.5\!\times\!10^{-3}$         & $2.2\!\times\!10^{-16}\,\star$ & $4.7\!\times\!10^{-2}$ \\
0.8 & $2.0\!\times\!10^{-15}\,\star$ & $2.2\!\times\!10^{-16}\,\star$ & $1.4\!\times\!10^{-2}$ \\
1.0 & $8.9\!\times\!10^{-16}\,\star$ & $2.8\!\times\!10^{-16}\,\star$ & $5.2\!\times\!10^{-2}$ \\
1.5 & $1.3\!\times\!10^{-3}$         & $3.3\!\times\!10^{-16}\,\star$ & $4.8\!\times\!10^{-2}$ \\
2.0 & $2.2\!\times\!10^{-16}\,\star$ & $2.2\!\times\!10^{-16}\,\star$ & $2.3\!\times\!10^{-2}$ \\
3.0 & $5.8\!\times\!10^{-16}\,\star$ & $2.2\!\times\!10^{-16}\,\star$ & $3.7\!\times\!10^{-2}$ \\
5.0 & $9.2\!\times\!10^{-16}\,\star$ & $2.2\!\times\!10^{-16}\,\star$ & $2.9\!\times\!10^{-2}$ \\
\bottomrule
\end{tabular}
\end{center}

\textbf{Verdict.}
\begin{itemize}
  \item \textbf{Lin-CDF kernel}: ALL 9 reach machine $\varepsilon$ — fast, robust,
        but smoothed.
  \item \textbf{Cheb strict}: 6 of 9 reach machine $\varepsilon$ — exact REE, but
        slow and $N_Q$ sweet-spot-dependent.
  \item \textbf{Lin-CDF strict}: NONE reach machine $\varepsilon$ — floors at
        $\sim 10^{-2}$. The piecewise-linear $P$ has piecewise-constant
        $\partial P/\partial u$, which generates step discontinuities in the
        co-area integrand that POU cannot smooth out.
\end{itemize}

\paragraph{Why strict $h{=}0$ needs the Chebyshev representation.}
The co-area integrand $f f / |\nabla P|$ depends pointwise on the derivative
$\nabla P$. On the linear-CDF grid the derivative is piecewise constant and
jumps at every grid line. As the root location $\xi_b^*(p)$ moves continuously
with $p$, the value of $|\nabla P|$ at the root changes discontinuously
each time the root crosses a grid line. POU smooths $|\nabla P| \to 0$
(turning points) but cannot smooth $|\nabla P|$ JUMPS. Chebyshev gives a
$C^\infty$ derivative everywhere, so no jumps.

\section{FP slopes and deficits}

\begin{center}
\includegraphics[width=\linewidth]{01_slope_three.png}
\end{center}

The three operators give three different FPs:

\begin{center}
\begin{tabular}{cccc}
\toprule
$\gamma$ & Cheb strict $\alpha^*$ & Lin-CDF kernel $\alpha^*$ & Lin-CDF strict $\alpha^*$ \\
\midrule
0.1 & 0.318 & 0.205 & 0.541 \\
0.3 & 0.304 & 0.228 & 0.550 \\
0.5 & 0.318 & 0.256 & 0.571 \\
0.8 & 0.329 & 0.284 & 0.652 \\
1.0 & 0.345 & 0.296 & 0.649 \\
1.5 & 0.368 & 0.314 & 0.658 \\
2.0 & 0.383 & 0.324 & 0.606 \\
3.0 & 0.415 & 0.336 & 0.711 \\
5.0 & 0.449 & 0.346 & 0.734 \\
\bottomrule
\end{tabular}
\end{center}

\begin{center}
\includegraphics[width=\linewidth]{03_deficit_three.png}
\end{center}

\paragraph{Pattern.} For each $\gamma$:
\begin{itemize}
  \item Lin-CDF kernel slope $<$ Cheb strict slope $<$ Lin-CDF strict slope
  \item The kernel smoothing biases the FP DOWN (less revealing).
  \item The strict-h=0-on-linear-grid biases UP (over-revealing) and never
        converges fully because of derivative jumps.
  \item Cheb strict is the only operator that's both exact AND converges
        to machine $\varepsilon$.
\end{itemize}

\paragraph{One-to-one breakdown.}
\begin{itemize}
  \item Cheb strict: $0.003-0.036$ (genuine breakdown captured).
  \item Lin-CDF kernel: $\sim 0$ for $\gamma \geq 0.5$ (smoothed away by
        kernel band).
  \item Lin-CDF strict: $0.001-0.008$ (intermediate, partly reflective
        of the unconverged FP).
\end{itemize}

\section{Contours at $\gamma=1$, $u_3=0$}

\begin{center}
\includegraphics[width=\linewidth]{05_contours_g1_three.png}
\end{center}

Same equilibrium concept, three different numerical pictures.

\section{Lin-CDF strict contour gallery}

\begin{center}
\includegraphics[width=\linewidth]{04_lin_strict_contours.png}
\end{center}

\section{logit-$P$ vs $T$ at $\gamma=1$}

\begin{center}
\includegraphics[width=\linewidth]{06_scatter_three_g1.png}
\end{center}

\section{Summary}

\begin{itemize}
  \item \textbf{For the actual REE} at strict $h{=}0$: use Cheb strict POU.
        Slope $\alpha^* = 0.345$ at $\gamma{=}1$, one-to-one deficit $0.0285$.
  \item \textbf{For fast parametric sweeps}: use Lin-CDF kernel band. Fast,
        robust, all $\gamma$ reach machine $\varepsilon$ — but FP is the
        regularised version (smaller slope, near-zero one-to-one deficit).
  \item \textbf{Lin-CDF strict $h{=}0$ does not work} for high-precision FP at this
        $G$ — derivative jumps prevent Newton from converging below
        $10^{-2}$.
\end{itemize}

The Chebyshev representation is \emph{required} for strict-$h{=}0$
machine-$\varepsilon$ convergence: it provides the $C^\infty$ derivatives
that the strict co-area integrand needs.

\end{document}
